% Appendix A -- The run-time library interface.  One entry per global entry
% point of the run-time objects (seed/rt/start.asm, start_dos.asm,
% fileio.asm) and per string routine Turbo Initech writes into a program
% (os/tps/tps.pas the __str_* prelude; os/tps/CODEGEN.md s5); KIND
% rtlentry of the generated inventory, both directions.  Sources (Law 1):
% the assembly sources named above (register use, INT 21h functions, the
% mutation entries SEED_MUT_FILEIO_SHORT_WRITE / _WRONG_HANDLE and the gate
% that uses them, test-seed-fileio-os-mutant).
\Gates{test-tps-gen-fpc, test-tps-gen-os, test-seed-fileio-fpc, test-seed-fileio-build, test-seed-fileio-os-mutant, test-compiler-os, test-pubs-tps}
\NoLetterTabs
\chapter{The run-time library interface}
\ix{run-time library, interface}
\ix{entry points}

This appendix describes, for programmers who read the assembly language
written by Turbo Initech, every entry point that a compiled program calls
or provides.  The entry points come from three places:

\begin{itemize}
\item the start object, linked first into every program, which contains
  the entry point \code{\_start} and the output routines;
\item the file object, linked into every program, which contains the
  \code{rtl\_file\_} routines;
\item the program itself: when a program uses strings, Turbo Initech writes
  the four \code{\_\_str\_} routines into its \file{TPSOUT.S}.
\end{itemize}

There are two start objects.  The InitechDOS start object makes a
\file{.COM} program; the stand-alone start object, used on the development
system to test programs that do not use files, is loaded at the first
megabyte of memory by the development system's loader, has its own 16K
stack, and ends by writing to I/O port F4h and halting.  Both contain the
same output routines, byte for byte.

\section{The start object}

\begin{RtlEntry}{_start}{\_start}
\LFunction The entry point of every Turbo Initech program.
\LCalling Not called: InitechDOS starts the program here, at its first byte.
\LRemarks In the InitechDOS start object, \code{\_start} sets every byte of
the \code{.bss} section to zero, loads the stack pointer with 7FFFCh, calls
\code{pas\_main}, and ends the program with INT 21h function 4Ch and
exit code 0.
\LSeeAlso \code{pas\_main} (Programmer's Guide, Chapter 11)
\end{RtlEntry}

\begin{RtlEntry}{serial_putc}{serial\_putc}
\LFunction Writes one character to COM1.
\LCalling AL = the character; \code{call serial\_putc}.
\LRemarks Writes AL to the data port of COM1 (3F8h).  Changes DX.
\std{write}{Write} calls it for a \rw{char} argument and for the line end
of \std{writeln}{Writeln}.
\LSeeAlso \code{serial\_puts}, \code{serial\_put\_int}
\end{RtlEntry}

\begin{RtlEntry}{serial_puts}{serial\_puts}
\LFunction Writes a string of bytes ended by a zero byte to COM1.
\LCalling EAX = the address of the first byte; \code{call serial\_puts}.
\LRemarks Writes each byte up to, not including, the first zero byte.
Preserves ESI; changes EAX and DX.  \std{write}{Write} calls it for a string
literal and for the texts \code{TRUE} and \code{FALSE}.
\LSeeAlso \code{serial\_putc}
\end{RtlEntry}

\begin{RtlEntry}{serial_put_int}{serial\_put\_int}
\LFunction Writes a signed integer to COM1 in decimal.
\LCalling EAX = the integer; \code{call serial\_put\_int}.
\LRemarks Writes a minus sign for a negative value, then the digits, with
no leading zeros or spaces.  Changes EAX, EBX, ECX and EDX; preserves ESI.
\LSeeAlso \code{serial\_putc}
\end{RtlEntry}

\section{The file object}
The file routines use the C calling convention of the compiled code
(Programmer's Guide, Chapter 11): the arguments are pushed last first, the
caller removes them, and EBX, ESI, EDI and EBP are preserved.  Each takes
the address of a file variable as its first argument.  When InitechDOS
reports an error, each routine writes \code{FILEIO-ERROR AX=}\textit{code}
to COM1 and ends the program with exit code 1.

\begin{RtlEntry}{rtl_file_assign}{rtl\_file\_assign}
\LFunction The code of \std{assign}{Assign}.
\LCalling \code{push} the address of the name (a string with its length
byte), \code{push} the address of the file variable, \code{call}.
\LRemarks Closes the file if its handle is that of an open file, copies the
name into the file variable as bytes ended by a zero byte, and marks the
file closed.
\LSeeAlso \code{rtl\_file\_reset}, \code{rtl\_file\_rewrite}
\end{RtlEntry}

\begin{RtlEntry}{rtl_file_reset}{rtl\_file\_reset}
\LFunction The code of \std{reset}{Reset}.
\LCalling \code{push} the record size, \code{push} the address of the file
variable, \code{call}.
\LRemarks Refuses a record size other than 1 (code 1).  Closes the file if
it is open, then opens the named file for reading with INT 21h function
3Dh and stores the handle.
\LSeeAlso \code{rtl\_file\_blockread}
\end{RtlEntry}

\begin{RtlEntry}{rtl_file_rewrite}{rtl\_file\_rewrite}
\LFunction The code of \std{rewrite}{Rewrite}.
\LCalling As for \code{rtl\_file\_reset}.
\LRemarks Refuses a record size other than 1.  Closes the file if it is
open, then creates the named file, or empties it, with INT 21h function
3Ch and normal attributes, and stores the handle.
\LSeeAlso \code{rtl\_file\_blockwrite}
\end{RtlEntry}

\begin{RtlEntry}{rtl_file_blockread}{rtl\_file\_blockread}
\LFunction The code of \std{blockread}{BlockRead}.
\LCalling \code{push} the address of the result variable, \code{push} the
count, \code{push} the address of the buffer, \code{push} the address of
the file variable, \code{call}.
\LRemarks Refuses a negative count (code 1).  Reads with INT 21h function
3Fh and stores the number of bytes read in the result variable.
\LSeeAlso \code{rtl\_file\_blockwrite}
\end{RtlEntry}

\begin{RtlEntry}{rtl_file_blockwrite}{rtl\_file\_blockwrite}
\LFunction The code of \std{blockwrite}{BlockWrite}.
\LCalling As for \code{rtl\_file\_blockread}.
\LRemarks Refuses a negative count.  Writes with INT 21h function 40h and
stores the number of bytes written in the result variable.
\LSeeAlso \code{rtl\_file\_blockread}
\end{RtlEntry}

\begin{RtlEntry}{rtl_file_blockwrite_short}{rtl\_file\_blockwrite\_short}
\LFunction A test entry that writes one byte less than asked and reports
the full count.
\LCalling As for \code{rtl\_file\_blockwrite}.
\LRemarks Turbo Initech never calls this entry.  It exists so that the
tests of the run-time library can prove that they detect a write that
loses data.
\LSeeAlso \code{rtl\_file\_blockwrite\_wrong\_handle}
\end{RtlEntry}

\begin{RtlEntry}{rtl_file_blockwrite_wrong_handle}{rtl\_file\_blockwrite\_wrong\_handle}
\LFunction A test entry that writes to the handle after the file's own.
\LCalling As for \code{rtl\_file\_blockwrite}.
\LRemarks Turbo Initech never calls this entry.  It exists so that the
tests of the run-time library can prove that they detect a write to the
wrong file.
\LSeeAlso \code{rtl\_file\_blockwrite\_short}
\end{RtlEntry}

\section{The string routines}
Turbo Initech writes these four routines into \file{TPSOUT.S} when, and
only when, the program uses strings.  They take their arguments in
registers.  A string is addressed by its length byte.

\begin{RtlEntry}{__str_assign}{\_\_str\_assign}
\LFunction Copies a string into a string variable, cutting it to the
variable's capacity.
\LCalling ESI = the source, EDI = the target, ECX = the capacity of the
target; \code{call \_\_str\_assign}.
\LRemarks Stores the smaller of the source length and the capacity as the
length of the target, then copies that many characters.
\LSeeAlso \code{\_\_str\_concat}
\end{RtlEntry}

\begin{RtlEntry}{__str_concat}{\_\_str\_concat}
\LFunction Appends one string to another.
\LCalling EDI = the string appended to, ESI = the string appended;
\code{call \_\_str\_concat}.
\LRemarks Appends as many characters as fit in 255 and sets the new length.
\LSeeAlso \code{\_\_str\_assign}
\end{RtlEntry}

\begin{RtlEntry}{__str_cmp}{\_\_str\_cmp}
\LFunction Compares two strings.
\LCalling ESI = the left string, EDI = the right string;
\code{call \_\_str\_cmp}.
\LRemarks Compares the characters as unsigned bytes up to the shorter
length, then the lengths.  Returns $-1$, 0 or 1 in EAX for less than, equal
to and greater than.
\LSeeAlso \code{\_\_str\_write}
\end{RtlEntry}

\begin{RtlEntry}{__str_write}{\_\_str\_write}
\LFunction Writes the characters of a string to COM1.
\LCalling ESI = the string; \code{call \_\_str\_write}.
\LRemarks Writes as many characters as the length byte gives, one at a time
through \code{serial\_putc}.
\LSeeAlso \code{serial\_putc}
\end{RtlEntry}
